\documentclass[11pt,a4paper]{article}
\usepackage[T1]{fontenc}
\usepackage[utf8]{inputenc}
\usepackage{lmodern}
\usepackage[english]{babel}
\usepackage{amsmath,amssymb,amsthm}
\usepackage{booktabs}
\usepackage{longtable}
\usepackage{hyperref}
\usepackage[margin=1in]{geometry}
\usepackage{enumitem}
\usepackage{xcolor}

\hypersetup{
  colorlinks=true,
  linkcolor=blue,
  urlcolor=blue,
  citecolor=blue
}

\setlist{nosep,leftmargin=1.5em}

\begin{document}

\title{\Huge\textbf{LLM Learning Plan \& Reading List}\\[0.5cm]
	\large Resources, Courses, and Instrumental Papers}
\author{Bernard}
\date{\today}
\maketitle

\section*{How to Use This Plan}

This plan is sequenced so that each step builds on the previous. The first three steps
build theoretical foundations; the middle steps bridge to engineering; the final steps
reach the research frontier. Within each step, the first resource is the primary
recommendation; the second is a complementary alternative or deeper dive.

\vspace{0.5cm}

% ============================================================
% LEARNING PLAN
% ============================================================
\section{Learning Plan}

\begin{enumerate}[leftmargin=2em,label=\textbf{Step \arabic*:},itemsep=0.7cm]

	% --- Step 1 ---
	\item \textbf{Information Theory}\hfill\textit{foundations --- 2--4 weeks}

	      \textbf{1.} Cover \& Thomas, \href{https://www.wiley.com/en-us/Elements+of+Information+Theory\%2C+2nd+Edition-p-9780471241959}{Elements of Information Theory}, 2nd ed., Wiley 2006.
	      \hfill\textit{Ch.1--4, 8--10 (skippable: Ch.5--7)}

	      \textbf{2.} MacKay, \href{https://www.inference.org.uk/mackay/itila/}{Information Theory, Inference, and Learning Algorithms}, CUP 2003. \hfill\textit{Free online}
	      \newline \qquad Unique: marries info theory to ML/Bayesian inference in one unified narrative.

	      % --- Step 2 ---
	\item \textbf{Deep Learning Mathematics}\hfill\textit{foundations --- 4--6 weeks}

	      \textbf{1.} Goodfellow et al., \href{https://www.deeplearningbook.org/}{Deep Learning}, MIT Press 2016. \hfill\textit{Free online}
	      \newline \qquad Ch.1--6 (math + ML basics), Ch.9--10 (CNNs, RNNs/attention).

	      \textbf{2.} Chen \& Siegel, \href{https://mathdl.github.io/}{Mathematical Foundations of Deep Learning Models and Algorithms}, AMS 2025.
	      \newline \qquad Rigorous proofs of UAT, NTK, mean-field, backward propagation. Ch.13 covers transformers.

	      % --- Step 3 ---
	\item \textbf{LLM Architecture \& Theory}\hfill\textit{core theory --- 3--5 weeks}

	      \textbf{1.} Vaswani et al., \href{https://arxiv.org/abs/1706.03762}{Attention Is All You Need}, NeurIPS 2017. \hfill\textit{Primary source}
	      \newline \qquad Read alongside \href{https://arxiv.org/abs/2601.22170}{Large Language Models: A Mathematical Formulation}
	      (2026) for a modern mathematical treatment covering tokenisation, training, RLHF.

	      \textbf{2.} \href{https://arxiv.org/abs/2602.22271}{Support Tokens, Stability Margins, and a New Foundation for Robust LLMs} (2026).
	      \newline \qquad Re-frames causal self-attention as a large-margin classifier; introduces the ``support tokens'' concept and a Bayesian MAP penalty.

	      % --- Step 4 ---
	\item \textbf{GPU Architecture \& CUDA}\hfill\textit{hardware --- 4--6 weeks}

	      \textbf{1.} \href{https://docs.nvidia.com/cuda/cuda-c-programming-guide/}{CUDA C++ Programming Guide} (NVIDIA official). \hfill\textit{Reference}
	      \newline \qquad Supplement with \href{https://github.com/siboehm/CUDAByExamples}{siboehm, CUDA by Examples} (blog + code).

	      \textbf{2.} \href{https://github.com/wafer-ai/gpu-perf-engineering-resources}{GPU Performance Engineering Resources}, wafer-ai.
	      \hfill\textit{Curated list}
	      \newline \qquad The definitive currated guide: roofline model, shared memory, tensor cores, flash attention, PagedAttention, CUDA graphs, warp specialisation.

	      % --- Step 5 ---
	\item \textbf{LLM Inference Optimisation}\hfill\textit{engineering --- 3--4 weeks}

	      \textbf{1.} \href{https://localaimaster.com/blog/cuda-optimization-local-llms}{CUDA Optimization Techniques for Local LLMs: The Complete 2026 Guide}.
	      \newline \qquad Practical ranking of optimisations by speedup/effort/risk: FlashAttention, KV-cache quantisation, FP8, speculative decoding, CUDA graphs.

	      \textbf{2.} \href{https://arxiv.org/abs/2302.01318}{Accelerating Large Language Model Decoding with Speculative Decoding}
	      (Chen et al., 2023) + \href{https://arxiv.org/abs/2205.14135}{FlashAttention} (Dao et al., 2022).
	      \newline \qquad Read both end-to-end; then study the KV cache in the \href{https://arxiv.org/abs/2309.06180}{PagedAttention} paper (Kwon et al., 2023).

	      % --- Step 6 ---
	\item \textbf{Production Serving}\hfill\textit{deployment --- 2--3 weeks}

	      \textbf{1.} \href{https://docs.vllm.ai/}{vLLM documentation} + \href{https://github.com/vllm-project/vllm}{GitHub}.
	      \newline \qquad Set up a local server; understand continuous batching, prefix caching, tensor parallelism, AWQ quantisation.

	      \textbf{2.} \href{https://github.com/NVIDIA/TensorRT-LLM}{TensorRT-LLM} (NVIDIA) + \href{https://github.com/ggerganov/llama.cpp}{llama.cpp}.
	      \newline \qquad Compare the two: TRT-LLM for max throughput on datacenter GPUs, llama.cpp for edge/CPU inference. Benchmark both on the same model.

	      % --- Step 7 ---
	\item \textbf{Software Engineering for LLMs}\hfill\textit{application --- 2--3 weeks}

	      \textbf{1.} Lewis et al., \href{https://arxiv.org/abs/2005.11401}{Retrieval-Augmented Generation for Knowledge-Intensive NLP Tasks}, NeurIPS 2020.
	      \newline \qquad The foundational RAG paper. Then build a RAG pipeline with \href{https://docs.llamaindex.ai/}{LlamaIndex}.
	      \newline \qquad Supplement: Wei et al., \href{https://arxiv.org/abs/2201.11903}{Chain-of-Thought Prompting Elicits Reasoning}, NeurIPS 2022.

	      \textbf{2.} Implement an LLM agent with \href{https://python.langchain.com/}{LangChain} that uses tools (calculator, search, database).
	      \newline \qquad Study: Ouyang et al., \href{https://arxiv.org/abs/2203.02155}{Training Language Models to Follow Instructions with Human Feedback}, NeurIPS 2022.

	      % --- Step 8 ---
	\item \textbf{Research Frontier: Long Context, Sparsity, RLVR}\hfill\textit{cutting edge --- 4--6 weeks}

	      \textbf{1.} \href{https://arxiv.org/abs/2603.22300}{Sparse Feature Attention: Dimension-Level Sparsity for Efficient Long-Context LLMs} (2026).
	      \newline \qquad Proposes learning $k$-sparse query/key codes; introduces FlashSFA. Matches dense perplexity at 2.5$\times$ speedup.

	      \textbf{2.} \href{https://arxiv.org/abs/2603.02146}{LongRLVR: Reinforcement Learning with Verifiable Rewards for Long-Context Reasoning} (2026).
	      \newline \qquad Proves outcome-only RLVR suffers vanishing gradients for contextual grounding; introduces a dense context reward. Improves RULER from 17~$\to$~88.

	      % --- Step 9 ---
	\item \textbf{Depth Scaling, Representation Geometry}\hfill\textit{emerging theory --- 3--4 weeks}

	      \textbf{1.} \href{https://openreview.net/pdf?id=UIQCQOQStf}{Why Gradual Depth Growth Improves Reasoning in LLMs} (2026, ICLR).
	      \newline \qquad Mechanistic analysis: depth-grown models use late layers; standard models suffer ``Curse of Depth'' (late layers are redundant). LIDAS algorithm.

	      \textbf{2.} \href{https://arxiv.org/pdf/2602.00217}{Understanding Embedding Condensation in LLMs} (2026).
	      \newline \qquad Smaller models exhibit collapsed embedding geometry; proposes dispersion loss that recovers the representational diversity of larger models.

\end{enumerate}

\newpage

% ============================================================
% INSTRUMENTAL PAPERS
% ============================================================
\section{Instrumental Papers}

\subsection*{Foundations}

\begin{enumerate}[leftmargin=2em,label=(\arabic*)]
	\item \textbf{Attention Is All You Need.} Vaswani et al. NeurIPS 2017.
	      \href{https://arxiv.org/abs/1706.03762}{arxiv.org/abs/1706.03762}
	      \newline The paper that started everything. Transformer architecture, multi-head attention, positional encodings.

	\item \textbf{Scaling Laws for Neural Language Models.} Kaplan et al. 2020.
	      \href{https://arxiv.org/abs/2001.08361}{arxiv.org/abs/2001.08361}
	      \newline Power-law scaling of loss with parameters, data, and compute. Established the empirical foundation.

	\item \textbf{Training Compute-Optimal Large Language Models.} Hoffmann et al. (Chinchilla). NeurIPS 2022.
	      \href{https://arxiv.org/abs/2203.15556}{arxiv.org/abs/2203.15556}
	      \newline Models and data should be scaled equally. Refuted the Kaplan scaling recommendations.

	\item \textbf{A Mathematical Theory of Communication.} Shannon. Bell System Technical Journal 1948.
	      \href{https://people.math.harvard.edu/~ctm/home/text/others/shannon/entropy/entropy.pdf}{Free PDF}
	      \newline The intellectual foundation of all information-theoretic reasoning about LLMs.

	\item \textbf{The Information Bottleneck Method.} Tishby, Pereira, Bialek. Allerton 1999.
	      \href{https://arxiv.org/abs/physics/0004057}{arxiv.org/abs/physics/0004057}
	      \newline Framework for understanding hidden representations as compressing input while preserving task signal.
\end{enumerate}

\subsection*{Architecture and Training}

\begin{enumerate}[leftmargin=2em,label=(\arabic*),start=6]
	\item \textbf{Language Models are Few-Shot Learners (GPT-3).} Brown et al. NeurIPS 2020.
	      \href{https://arxiv.org/abs/2005.14165}{arxiv.org/abs/2005.14165}
	      \newline Demonstrated in-context learning at scale; the paradigm shift from fine-tuning to prompting.

	\item \textbf{Training Language Models to Follow Instructions with Human Feedback (InstructGPT).}
	      Ouyang et al. NeurIPS 2022. \href{https://arxiv.org/abs/2203.02155}{arxiv.org/abs/2203.02155}
	      \newline The canonical RLHF paper: PPO + KL penalty for alignment.

	\item \textbf{Llama 2: Open Foundation and Fine-Tuned Chat Models.} Touvron et al. 2023.
	      \href{https://arxiv.org/abs/2307.09288}{arxiv.org/abs/2307.09288}
	      \newline GQA, SwiGLU, RoPE, RLHF recipe. The reference open-source architecture.

	\item \textbf{The Llama 3 Herd of Models.} Dubey et al. 2024.
	      \href{https://arxiv.org/abs/2407.21783}{arxiv.org/abs/2407.21783}
	      \newline Scaling to 405B; instruction tuning, tool use, multilingual, long-context.

	\item \textbf{DeepSeek-R1: Incentivizing Reasoning Capability in LLMs via Reinforcement Learning.}
	      Guo et al. 2025. \href{https://arxiv.org/abs/2501.12948}{arxiv.org/abs/2501.12948}
	      \newline RLVR with verifiable rewards produces emergent reasoning (the ``aha moment'').
\end{enumerate}

\subsection*{Inference and Systems}

\begin{enumerate}[leftmargin=2em,label=(\arabic*),start=11]
	\item \textbf{FlashAttention: Fast and Memory-Efficient Exact Attention with IO-Awareness.}
	      Dao et al. NeurIPS 2022. \href{https://arxiv.org/abs/2205.14135}{arxiv.org/abs/2205.14135}
	      \newline Tiling + recomputation eliminates $O(N^2)$ memory. Essential reading for GPU optimisation.

	\item \textbf{Efficient Memory Management for Large Language Model Serving with PagedAttention.}
	      Kwon et al. (vLLM). SOSP 2023. \href{https://arxiv.org/abs/2309.06180}{arxiv.org/abs/2309.06180}
	      \newline Virtual memory for KV cache; eliminates fragmentation. Continuous batching.

	\item \textbf{Fast Inference from Transformers via Speculative Decoding.}
	      Leviathan et al. ICML 2023. \href{https://arxiv.org/abs/2211.17192}{arxiv.org/abs/2211.17192}
	      \newline Draft-verify paradigm; expected speedup analysis.

	\item \textbf{GPTQ: Accurate Post-Training Quantization for Generative Pre-Trained Transformers.}
	      Frantar et al. ICLR 2023. \href{https://arxiv.org/abs/2210.17323}{arxiv.org/abs/2210.17323}
	      \newline Optimal Brain Quantiser adapted for LLMs; the standard PTQ baseline.

	\item \textbf{AWQ: Activation-aware Weight Quantization for LLM Compression and Acceleration.}
	      Lin et al. MLSys 2024. \href{https://arxiv.org/abs/2306.00978}{arxiv.org/abs/2306.00978}
	      \newline Observes 1\% of weights (salient channels) are disproportionately important; protects them during quantisation.
\end{enumerate}

\subsection*{Engineering and Methods}

\begin{enumerate}[leftmargin=2em,label=(\arabic*),start=16]
	\item \textbf{Chain-of-Thought Prompting Elicits Reasoning in Large Language Models.}
	      Wei et al. NeurIPS 2022. \href{https://arxiv.org/abs/2201.11903}{arxiv.org/abs/2201.11903}
	      \newline CoT increases effective depth; enables problems unsolvable by direct prompting.

	\item \textbf{Retrieval-Augmented Generation for Knowledge-Intensive NLP Tasks.}
	      Lewis et al. NeurIPS 2020. \href{https://arxiv.org/abs/2005.11401}{arxiv.org/abs/2005.11401}
	      \newline The RAG paradigm: index, retrieve, generate.

	\item \textbf{Tree of Thoughts: Deliberate Problem Solving with Large Language Models.}
	      Yao et al. NeurIPS 2023. \href{https://arxiv.org/abs/2305.10601}{arxiv.org/abs/2305.10601}
	      \newline Search over reasoning trees; connects LLM decoding to classical AI search.

	\item \textbf{HarmonyDream: No Need for Task-Specific Decoders in World Models.}
	      Robine et al. ICLR 2025. \href{https://arxiv.org/abs/2502.12872}{arxiv.org/abs/2502.12872}
	      \newline Architectural innovation for latent-state world models; demonstrates shared representations.

	\item \textbf{React: Synergizing Reasoning and Acting in Language Models.}
	      Yao et al. ICLR 2023. \href{https://arxiv.org/abs/2210.03629}{arxiv.org/abs/2210.03629}
	      \newline Interleaved reasoning traces and tool actions; the foundation of modern agent frameworks.
\end{enumerate}

\subsection*{Research Frontier (2025--2026)}

\begin{enumerate}[leftmargin=2em,label=(\arabic*),start=21]
	\item \textbf{Large Language Models: A Mathematical Formulation.} Brennan et al. 2026.
	      \href{https://arxiv.org/abs/2601.22170}{arxiv.org/abs/2601.22170}
	      \newline Rigorous mathematical framework covering tokenisation, architecture, training, and deployment.

	\item \textbf{Support Tokens, Stability Margins, and a New Foundation for Robust LLMs.} 2026.
	      \href{https://arxiv.org/abs/2602.22271}{arxiv.org/abs/2602.22271}
	      \newline Causal self-attention induces a margin-to-degeneracy barrier; support tokens emerge analogous to SVMs.

	\item \textbf{Language Models are Almost Surely Injective.} 2026.
	      \href{https://arxiv.org/abs/2510.15511}{arxiv.org/abs/2510.15511}
	      \newline Proves decoder-only transformers are injective; SIPIT algorithm recovers input from hidden states.

	\item \textbf{Sparse Feature Attention.} 2026.
	      \href{https://arxiv.org/abs/2603.22300}{arxiv.org/abs/2603.22300}
	      \newline Dimension-level sparsity for attention; FlashSFA kernel. $k$-sparse queries/keys. 2.5$\times$ speedup.

	\item \textbf{LongRLVR: RL with Verifiable Rewards for Long-Context Reasoning.} 2026.
	      \href{https://arxiv.org/abs/2603.02146}{arxiv.org/abs/2603.02146}
	      \newline Proves outcome-only RLVR has vanishing gradients for contextual grounding. Dense context reward.

	\item \textbf{Why Gradual Depth Growth Improves Reasoning in LLMs.} 2026 (ICLR).
	      \href{https://openreview.net/pdf?id=UIQCQOQStf}{OpenReview}
	      \newline Depth-grown models escape the Curse of Depth; late layers remain useful. LIDAS algorithm.

	\item \textbf{Understanding Embedding Condensation in LLMs.} 2026.
	      \href{https://arxiv.org/pdf/2602.00217}{arxiv.org/abs/2602.00217}
	      \newline Smaller models collapse embedding diversity; dispersion loss recovers large-model geometry.

	\item \textbf{In-Context Policy Optimization.} 2026 (ICLR).
	      \href{https://arxiv.org/abs/2603.01335}{arxiv.org/abs/2603.01335}
	      \newline Theoretical framework proving one-layer linear self-attention can execute policy optimisation in-context.
\end{enumerate}

\vfill
\begin{center}
	\rule{0.5\textwidth}{0.5pt}\\[0.3cm]
	\textit{Last updated: \today}
\end{center}

\end{document}
